% GENERATED FILE. Edit canonical lessons, then run npm run generate:books.
% canonical-source-hash: fnv1a64:50cae122e12f9803
% canonical-lessons: MR-C390-ekun, MR-C390-ekandarit, MR-C390-tattva, MR-C390-vises, MR-C390-vyatirikta

\chapter{On the Whole}
\label{ch:mr-on-the-whole}
\hlchaptermodality{390}

\begin{chapteropening}
\textbf{By the end of this chapter:} \emph{I can sum up with in all, on the whole and apart from.}

Complete the last lesson of chapter 390: I can sum up with in all, on the whole and apart from.
\end{chapteropening}

\section[\texorpdfstring{\mr{एकूण} (ekūṇ)}{एकूण (ekūṇ)}]{\mr{एकूण} (ekūṇ) — in total; overall}

\label{lesson:MR-C390-ekun}



\begin{quote}
Before the new one: say the Marathi for purely, then the Marathi for usually.
\end{quote}

\subsection*{You'll want to know: \mr{एकूण}}
\textbf{\mr{एकूण}} — \emph{ekūṇ} — \textquotedblleft{}in total; overall\textquotedblright{}.

\begin{quote}
\mr{वर्गात} \mr{एकूण} \mr{तीस} \mr{मुलं} \mr{आहेत}.
\end{quote}

\begin{quote}
\emph{vargāt ekūṇ tīs mulaṁ āhet.}
\end{quote}

\begin{quote}
There are thirty children in the class in all.
\end{quote}

Also an adjective: \mr{एकूण} \mr{खर्च}, the total cost.

\begin{cousinweb}
\textbf{\mr{एकूण}} (\emph{ekūṇ}) is a Marathi word whose origin is not settled.
\end{cousinweb}

\begin{grammarlens}[title={What you've built}]
One more word for saying what happens if.
\end{grammarlens}

\subsection*{Your turn}
\begin{itemize}
\raggedright
  \item \emph{Say it:} \emph{ekūṇ}
  \item \emph{Say it:} \emph{ekūṇ}, once more
  \item \emph{Say it:} \emph{vargāt ekūṇ tīs mulaṁ āhet.} — There are thirty children in the class in all.
  \item \emph{Recall:} say the Marathi for purely, then the Marathi for usually, then say \emph{ekūṇ} again
\end{itemize}

\begin{tcolorbox}[breakable,colback=teal!4,colframe=teal!35!black,title={Before you move on}]
What does \mr{एकूण} mean? (\textquotedblleft{}In total; overall\textquotedblright{}.) Say it once more.
\end{tcolorbox}
\section[\texorpdfstring{\mr{एकंदरीत} (ekandarīt)}{एकंदरीत (ekandarīt)}]{\mr{एकंदरीत} (ekandarīt) — on the whole}

\label{lesson:MR-C390-ekandarit}



\begin{quote}
Before the new one: say the Marathi for usually, then the Marathi for in total.
\end{quote}

\subsection*{You'll want to know: \mr{एकंदरीत}}
\textbf{\mr{एकंदरीत}} — \emph{ekandarīt} — \textquotedblleft{}on the whole\textquotedblright{}.

\begin{quote}
\mr{एकंदरीत} \mr{सहल} \mr{छान} \mr{झाली}.
\end{quote}

\begin{quote}
\emph{ekandarīt sahal chān jhālī.}
\end{quote}

\begin{quote}
On the whole, the trip went well.
\end{quote}

\mr{एकंदर}, whole, + -\mr{ईत}.

\begin{cousinweb}
\textbf{\mr{एकंदरीत}} (\emph{ekandarīt}) is built on \mr{एकंदर}, \textquotedblleft{}whole, total\textquotedblright{}, with the ending -\mr{ईत}, \textquotedblleft{}in\textquotedblright{}; the origin of \mr{एकंदर} is not settled.
\end{cousinweb}

\begin{grammarlens}[title={What you've built}]
One more word for saying what happens if.
\end{grammarlens}

\subsection*{Your turn}
\begin{itemize}
\raggedright
  \item \emph{Say it:} \emph{ekandarīt}
  \item \emph{Say it:} \emph{ekandarīt}, once more
  \item \emph{Say it:} \emph{ekandarīt sahal chān jhālī.} — On the whole, the trip went well.
  \item \emph{Recall:} say the Marathi for usually, then the Marathi for in total, then say \emph{ekandarīt} again
\end{itemize}

\begin{tcolorbox}[breakable,colback=teal!4,colframe=teal!35!black,title={Before you move on}]
What does \mr{एकंदरीत} mean? (\textquotedblleft{}On the whole\textquotedblright{}.) Say it once more.
\end{tcolorbox}
\section[\texorpdfstring{\mr{तत्त्व} (tattva)}{तत्त्व (tattva)}]{\mr{तत्त्व} (tattva) — a principle}

\label{lesson:MR-C390-tattva}



\begin{quote}
Before the new one: say the Marathi for in total, then the Marathi for on the whole.
\end{quote}

\subsection*{You'll want to know: \mr{तत्त्व}}
\textbf{\mr{तत्त्व}} — \emph{tattva} — \textquotedblleft{}a principle\textquotedblright{}.

\begin{quote}
\mr{खोटं} \mr{न} \mr{बोलणं} \mr{हे} \mr{त्याचं} \mr{तत्त्व} \mr{आहे}.
\end{quote}

\begin{quote}
\emph{khoṭaṁ na bolṇaṁ he tyācaṁ tattva āhe.}
\end{quote}

\begin{quote}
Not telling lies is his principle.
\end{quote}

Neuter.

\mr{तत्त्वतः}, in principle.

\begin{cousinweb}
\textbf{\mr{तत्त्व}} (\emph{tattva}) is a tatsama, taken from Sanskrit \emph{tattva}, \textquotedblleft{}real state, truth\textquotedblright{} (CDIAL 5642), from \emph{tat}, \textquotedblleft{}that\textquotedblright{}: thatness.
\end{cousinweb}

\begin{grammarlens}[title={What you've built}]
One more word for saying what happens if.
\end{grammarlens}

\subsection*{Your turn}
\begin{itemize}
\raggedright
  \item \emph{Say it:} \emph{tattva}
  \item \emph{Say it:} \emph{tattva}, once more
  \item \emph{Say it:} \emph{khoṭaṁ na bolṇaṁ he tyācaṁ tattva āhe.} — Not telling lies is his principle.
  \item \emph{Recall:} say the Marathi for in total, then the Marathi for on the whole, then say \emph{tattva} again
\end{itemize}

\begin{tcolorbox}[breakable,colback=teal!4,colframe=teal!35!black,title={Before you move on}]
What does \mr{तत्त्व} mean? (\textquotedblleft{}A principle\textquotedblright{}.) Say it once more.
\end{tcolorbox}
\section[\texorpdfstring{\mr{विशेष} (viśeṣ)}{विशेष (viśeṣ)}]{\mr{विशेष} (viśeṣ) — particular; special}

\label{lesson:MR-C390-vises}



\begin{quote}
Before the new one: say the Marathi for on the whole, then the Marathi for a principle.
\end{quote}

\subsection*{You'll want to know: \mr{विशेष}}
\textbf{\mr{विशेष}} — \emph{viśeṣ} — \textquotedblleft{}particular; special\textquotedblright{}.

\begin{quote}
\mr{आज} \mr{जेवणात} \mr{काही} \mr{विशेष} \mr{आहे} \mr{का}?
\end{quote}

\begin{quote}
\emph{āj jevṇāt kāhī viśeṣ āhe kā?}
\end{quote}

\begin{quote}
Is there anything special for dinner today?
\end{quote}

\mr{विशेष} \mr{करून}, especially.

\begin{cousinweb}
\textbf{\mr{विशेष}} (\emph{viśeṣ}) is a tatsama, taken from Sanskrit \emph{viśeṣa}, \textquotedblleft{}distinction, peculiarity\textquotedblright{} (Monier-Williams).
\end{cousinweb}

\begin{grammarlens}[title={What you've built}]
One more word for saying what happens if.
\end{grammarlens}

\subsection*{Your turn}
\begin{itemize}
\raggedright
  \item \emph{Say it:} \emph{viśeṣ}
  \item \emph{Say it:} \emph{viśeṣ}, once more
  \item \emph{Say it:} \emph{āj jevṇāt kāhī viśeṣ āhe kā?} — Is there anything special for dinner today?
  \item \emph{Recall:} say the Marathi for on the whole, then the Marathi for a principle, then say \emph{viśeṣ} again
\end{itemize}

\begin{tcolorbox}[breakable,colback=teal!4,colframe=teal!35!black,title={Before you move on}]
What does \mr{विशेष} mean? (\textquotedblleft{}Particular; special\textquotedblright{}.) Say it once more.
\end{tcolorbox}
\section[\texorpdfstring{\mr{व्यतिरिक्त} (vyatirikta)}{व्यतिरिक्त (vyatirikta)}]{\mr{व्यतिरिक्त} (vyatirikta) — apart from; besides}

\label{lesson:MR-C390-vyatirikta}



\begin{quote}
Before the new one: say the Marathi for a principle, then the Marathi for particular.
\end{quote}

\subsection*{You'll want to know: \mr{व्यतिरिक्त}}
\textbf{\mr{व्यतिरिक्त}} — \emph{vyatirikta} — \textquotedblleft{}apart from; besides\textquotedblright{}.

\begin{quote}
\mr{मराठीच्या} \mr{व्यतिरिक्त} \mr{त्याला} \mr{आणखी} \mr{दोन} \mr{भाषा} \mr{येतात}.
\end{quote}

\begin{quote}
\emph{marāṭhīcyā vyatirikta tyālā āṇkhī don bhāṣā yetāt.}
\end{quote}

\begin{quote}
Apart from Marathi he knows two more languages.
\end{quote}

After -\mr{च्या}, or attached: \mr{मराठीव्यतिरिक्त}.

\begin{cousinweb}
\textbf{\mr{व्यतिरिक्त}} (\emph{vyatirikta}) is a tatsama, taken from Sanskrit \emph{vyatirikta}, \textquotedblleft{}left apart, excluded\textquotedblright{} (Monier-Williams).
\end{cousinweb}

\begin{grammarlens}[title={What you've built}]
One more word for saying what happens if.
\end{grammarlens}

\subsection*{Your turn}
\begin{itemize}
\raggedright
  \item \emph{Say it:} \emph{vyatirikta}
  \item \emph{Say it:} \emph{vyatirikta}, once more
  \item \emph{Say it:} \emph{marāṭhīcyā vyatirikta tyālā āṇkhī don bhāṣā yetāt.} — Apart from Marathi he knows two more languages.
  \item \emph{Recall:} say the Marathi for a principle, then the Marathi for particular, then say \emph{vyatirikta} again
\end{itemize}

\begin{tcolorbox}[breakable,colback=teal!4,colframe=teal!35!black,title={Before you move on}]
What does \mr{व्यतिरिक्त} mean? (\textquotedblleft{}Apart from; besides\textquotedblright{}.) Say it once more.
\end{tcolorbox}
